%% LyX 2.4.4 created this file.  For more info, see https://www.lyx.org/.
%% Do not edit unless you really know what you are doing.
\documentclass[english]{article}
\usepackage[T1]{fontenc}
\usepackage[latin9]{inputenc}
\usepackage{enumitem}
\usepackage{amsmath}
\usepackage{amsthm}
\usepackage{cancel}
\usepackage{geometry}
\geometry{verbose,tmargin=1in,bmargin=1in,lmargin=1in,rmargin=1in}

\makeatletter

%%%%%%%%%%%%%%%%%%%%%%%%%%%%%% LyX specific LaTeX commands.
%% Because html converters don't know tabularnewline
\providecommand{\tabularnewline}{\\}

%%%%%%%%%%%%%%%%%%%%%%%%%%%%%% Textclass specific LaTeX commands.
\numberwithin{equation}{section}
\numberwithin{figure}{section}
\newlength{\lyxlabelwidth}      % auxiliary length 

%%%%%%%%%%%%%%%%%%%%%%%%%%%%%% User specified LaTeX commands.
\usepackage{fancyhdr}
\usepackage{lastpage}
\pagestyle{fancy}
\usepackage{enumitem}
\usepackage{ifthen}
%sets math fonts to be more readable
\everymath=\expandafter{\the\everymath\displaystyle}
%\DeclareMathSizes{10}{10}{10}{10}
\usepackage{multicol}

\usepackage{tikz}
\usetikzlibrary{plotmarks}
\usepackage{pgfplots}
\pgfplotsset{compat=newest}

\usepackage{calc}
\usepackage[nomessages]{fp}

\usepackage{advdate}    % Advancing/saving dates
\usepackage{datetime}   % Dates formatting
\usepackage{datenumber} % Counters for dates
% Added by lyx2lyx
\setlength{\parskip}{\smallskipamount}
\setlength{\parindent}{0pt}

\makeatother

\usepackage{babel}
\begin{document}
\renewcommand{\author}{Dr.\,James Ashe}

\newcommand{\classNum}{132}

\newcommand{\classSec}{C and K}

\newcommand{\nameType}{Name:}

\newcommand{\examType}{Worksheet 1.2 and 1.3}

\renewcommand{\date}{\formatdate{28}{1}{2014}}

%%First page's header%%%%%%%%%%%%%%%%%%%%%%
\fancypagestyle{firststyle} %%header settings for first pagr
{
	\fancyhead[L]{MTH \classNum{}(\classSec{}) \author{}\\
	\textbf{\examType{}} \date{}}
	\fancyhead[C]{\textbf{\nameType}}
	\setlength{\headsep}{14pt}
}
%%%%%%%%%%%%%%%%%%%%%%%%%%%%%%%%%%%%%%%%%%

%%Next page's header%%%%%%%%%%%%%%%%%%%%%%%
\setlist{nolistsep}
\lhead{\classNum{}(MTH \classSec{})} 
\chead{\textbf{\nameType}} 
\cfoot{\thepage{}~of~\pageref{LastPage}} 
\setlength{\headsep}{5pt} 
%%%%%%%%%%%%%%%%%%%%%%%%%%%%%%%%%%%%%%%%%%%

%%Various settings; see the preamble for other settings%%%%%
%%\setenumerate[0]{leftmargin=12pt,itemindent=0pt,labelwidth=0pt}
\renewcommand{\labelenumi}{\textbf{\arabic{enumi}.}}
\renewcommand{\labelenumii}{\textbf{\alph{enumii}.}}
\thispagestyle{firststyle}
%%%%%%%%%%%%%%%%%%%%%%%%%%%%%%%%%%%%%%%%%%
\newcommand{\xmax}{10}
\newcommand{\xmin}{-10}
\newcommand{\xstep}{0} %must be xmin+n
\newcommand{\ymax}{10}
\newcommand{\ymin}{-10}
\newcommand{\ystep}{0} %must be ymin+n

\newcommand{\xspace}{\vspace{1.2cm}}

\textbf{For exercises 1-2, let $f(x)=7-5x$ and $g(x)=2x-3$. Find
the following.}
\begin{enumerate}
\item $f(2)$\xspace
\item $g(1.5)$\xspace
\end{enumerate}
\textbf{Solve the problem.}
\begin{enumerate}[resume]
\item Let the supply and demand functions for butter pecan ice cream be
given by 
\[
p=S(q)=\frac{2}{5}q\mbox{ and }p=D(q)=100-\frac{2}{5}q\mbox{,}
\]
where $p$ is the price in dollars and $q$ is the number of 10-gallon
tubs. Find the equilibrium quantity and the equilibrium price.\vfill
\item FedUps delivers packages which cost \$2.00 per package to deliver.
The fixed cost to run the delivery truck is \$324 per day. If the
company charges \$5.00 per package, how many packages must be delivered
to make a profit of \$42?\vfill
\end{enumerate}
\textbf{SHORT ANSWER. Write the word or phrase that best completes
each statement or answers the question.}
\begin{enumerate}[resume]
\item In one or more paragraph, explain what it means to ``break even''
in business.\xspace\newpage
\end{enumerate}
\textbf{Find the equation of the least squares line.}
\begin{enumerate}[resume]
\item \label{enu:The-paired-data}The paired data below consists of the
average pupil-teacher ratio in public schools for 5 randomly selected
years.\\
\begin{tabular}{c|c|c|c|c|c}
Year & 1990 & 1994 & 1998 & 2002 & 2006\tabularnewline
\hline 
Ratio & 17.4 & 17.7 & 16.9 & 16.2 & 15.8\tabularnewline
\end{tabular}\vfill\xspace\xspace
\item \label{enu:The-paired-data-1}The paired data below consists of customers
(in millions) for 5 selected years.\\
\begin{tabular}{c|c|c|c|c|c}
Year & 2000 & 2002 & 2004 & 2006 & 2008\tabularnewline
\hline 
Customers  & 8.5 & 19.3 & 25.4 & 32.6 & 40.4\tabularnewline
\end{tabular}\vfill\xspace\xspace
\end{enumerate}
\textbf{SHORT ANSWER. Write the word or phrase that best completes
each statement or answers the question.}
\begin{enumerate}[resume]
\item Which sets of paired data from problems \ref{enu:The-paired-data}
and \ref{enu:The-paired-data-1} have a stronger linear relation?
(i.e., of the two equations which gives a better prediction for future
values?)\vfill
\end{enumerate}

\end{document}
